\chapter{Layered Architecture}
\label{bab:lapisan}

\section*{Tujuan Pembelajaran}

\begin{enumerate}
  \item Menjelaskan pembagian tanggung jawab pada keempat lapisan \textit{layered
    architecture} dan arah ketergantungan yang berlaku di antaranya.
  \item Membedakan lapisan tertutup dari lapisan terbuka, lalu menentukan kapan
    masing-masing bentuk layak dipilih.
  \item Mengukur jumlah fungsi penerus, pelanggaran arah ketergantungan, dan
    selisih waktu antara kedua bentuk pada satu aplikasi yang sama.
\end{enumerate}

\section{Pendahuluan}

\textit{Layered architecture} menata sebuah aplikasi menjadi beberapa lapisan
yang bertumpuk. Setiap lapisan mengerjakan satu jenis urusan teknis, lalu
memanggil lapisan di bawahnya bila membutuhkan sesuatu. Bentuk ini juga dikenal
sebagai arsitektur \textit{n-tier}~\cite{microsoft2024ntier}. Bentuk ini menjadi pilihan pertama bagi
sebagian besar sistem, dan sering dipakai tanpa pertimbangan panjang ketika tim
belum memutuskan arsitektur lain.

Fowler~\cite{fowler2015layering} mencatat bahwa pembagian menjadi tampilan,
domain, dan data merupakan pembagian yang paling sering dipakai, sekaligus yang
paling sering dipakai tanpa alasan yang dipikirkan.

Salah satu alasan pemakaiannya bersifat organisasi, bukan teknis. Pembagian
lapisan mengikuti pembagian kerja yang sudah ada di banyak tim pengembang. Ada
yang mengerjakan antarmuka, ada yang mengerjakan aturan bisnis, dan ada yang
mengerjakan basis data. Struktur tim tersebut memetakan diri dengan rapi ke
struktur lapisan, sehingga setiap orang tahu berkas mana yang menjadi
tanggung jawabnya.

Kemudahan itu ada harganya. Ketika aplikasi membesar, pemisahan yang tadinya
menolong justru memperlambat perubahan. Satu fitur baru sering menyentuh
keempat lapisan sekaligus. Bab ini membahas pembagian lapisan, aturan
ketergantungan di antaranya, serta pilihan antara lapisan tertutup dan lapisan
terbuka. Pembahasan mengikuti kerangka yang disusun oleh Richards dan
Ford~\cite{richards2020fundamentals}.

\section{Pembagian Lapisan}

Bentuk baku \textit{layered architecture} terdiri atas empat lapisan.
Gambar~\ref{fig:lapisan-empat} menunjukkan susunannya beserta arah panggilan
yang diizinkan. Panah hanya menunjuk ke bawah, dan tidak ada panah yang
melompati satu lapisan.

\begin{figure}[htbp]
  \centering
  \begin{tikzpicture}[
    lapis/.style={draw, rounded corners=2pt, minimum width=0.62\textwidth,
      minimum height=1.05cm, align=center, font=\scriptsize},
    panah/.style={-{Stealth[length=2mm]}, thick}
  ]
    \node[lapis, fill=blue!8] (pres) {\textbf{Lapisan Presentasi}\\
      {\tiny\mdseries Menerima masukan, menampilkan hasil}};
    \node[lapis, fill=blue!8, below=0.55cm of pres] (bis) {\textbf{Lapisan Bisnis}\\
      {\tiny\mdseries Memeriksa dan menjalankan aturan bisnis}};
    \node[lapis, fill=blue!8, below=0.55cm of bis] (per) {\textbf{Lapisan Persistensi}\\
      {\tiny\mdseries Menerjemahkan data simpanan menjadi objek}};
    \node[lapis, fill=green!10, below=0.55cm of per] (bas) {\textbf{Lapisan Basis Data}\\
      {\tiny\mdseries Menyimpan data secara permanen}};

    \draw[panah] (pres) -- (bis);
    \draw[panah] (bis) -- (per);
    \draw[panah] (per) -- (bas);
  \end{tikzpicture}
  \caption{Empat lapisan baku beserta arah panggilan yang diizinkan}
  \label{fig:lapisan-empat}
\end{figure}

\textbf{Lapisan presentasi} menangani segala hal yang dilihat pengguna. Lapisan
ini menerima masukan, memanggil lapisan bisnis, lalu menampilkan hasilnya.
Lapisan ini tidak mengetahui asal data maupun cara data disimpan. Contohnya
menyusun teks \texttt{100 USD = 1.625.000 IDR} lengkap dengan pemisah ribuan,
menampilkan pesan galat ketika nominal ditolak, dan mengubah daftar kurs menjadi
baris-baris teks. Bila layar diganti menjadi API, seluruh isi lapisan ini
dibuang dan ditulis ulang.

\textbf{Lapisan bisnis} memuat aturan yang membuat sistem tersebut bernilai.
Pemeriksaan kelayakan, perhitungan, dan keputusan proses berada di sini.
Lapisan ini tidak mengetahui bentuk antarmuka maupun bentuk tabel basis data.
Contohnya menolak nominal yang bernilai nol atau negatif, menolak nominal yang
melewati batas wajar, memeriksa token pemanggil, dan memeriksa apakah peran
pemanggil berhak membaca kurs jual. Seluruh aturan tersebut tetap berlaku
walaupun layarnya diganti.

\textbf{Lapisan persistensi} menjembatani objek di dalam program dengan baris di
dalam basis data. Kueri Structured Query Language (SQL) ditulis di lapisan ini,
sehingga lapisan bisnis tidak perlu mengetahui nama tabel. Contohnya menjalankan
kueri pencarian satu kurs, mengubah baris hasil kueri menjadi tupel siap pakai,
dan mengembalikan nilai kosong ketika pasangan mata uang tidak ditemukan. Bila
basis datanya berganti, hanya isi lapisan inilah yang ditulis ulang.

\textbf{Lapisan basis data} menyimpan data secara permanen. Lapisan ini biasanya
berupa perangkat lunak tersendiri, bukan kode yang ditulis tim pengembang.
Contohnya tabel \texttt{kurs} beserta kolom \texttt{nilai} dan
\texttt{margin}, kunci utama berupa pasangan kode mata uang, serta aturan bahwa
satu pasangan hanya boleh muncul satu kali. Keputusan menyimpan margin pada
kolom terpisah diambil di lapisan ini, dan keputusan itulah yang memungkinkan
kurs referensi dibaca tanpa ikut membuka marginnya.

Pembagian tersebut memberi satu keuntungan yang jelas. Perubahan pada satu
urusan teknis tertahan di satu lapisan. Penggantian basis data, misalnya, cukup
menyentuh lapisan persistensi selama bentuk objek yang dikembalikannya tetap.
Tabel~\ref{tab:tanggung-jawab} merangkum tanggung jawab dan larangan setiap
lapisan.

\begin{table}[htbp]
  \centering
  \caption{Tanggung jawab dan larangan setiap lapisan}
  \label{tab:tanggung-jawab}
  \begin{tabularx}{\textwidth}{|l|X|X|}
    \hline
    \textbf{Lapisan} & \textbf{Tanggung jawab} & \textbf{Tidak boleh mengetahui} \\
    \hline
    Presentasi & Menerima masukan dan menampilkan hasil & Nama tabel, kueri, sumber data \\
    \hline
    Bisnis & Memeriksa aturan dan menghitung hasil & Bentuk antarmuka, bentuk tabel \\
    \hline
    Persistensi & Menerjemahkan baris menjadi objek & Aturan bisnis, kebutuhan tampilan \\
    \hline
    Basis Data & Menyimpan data secara permanen & Seluruh lapisan di atasnya \\
    \hline
  \end{tabularx}
\end{table}

Agar pembagian tersebut tidak berhenti sebagai daftar, Tabel~\ref{tab:contoh-lapisan}
menunjukkan isi setiap lapisan pada aplikasi kurs yang dipakai sepanjang bab
ini. Perhatikan bahwa satu kalimat kebutuhan biasanya jatuh ke satu lapisan
saja.

\begin{table}[htbp]
  \centering
  \caption{Contoh isi setiap lapisan pada aplikasi kurs}
  \label{tab:contoh-lapisan}
  \begin{tabularx}{\textwidth}{|l|X|X|}
    \hline
    \textbf{Lapisan} & \textbf{Contoh isinya} & \textbf{Kalimat kebutuhan} \\
    \hline
    Presentasi & Menyusun teks \texttt{100 USD = 1.625.000 IDR} & Hasil ditampilkan dengan pemisah ribuan \\
    \hline
    Bisnis & Menolak nominal nol, memeriksa token dan peran & Hanya staf boleh melihat kurs jual \\
    \hline
    Persistensi & Mengubah baris tabel menjadi tupel & Kurs disimpan per pasangan mata uang \\
    \hline
    Basis Data & Tabel \texttt{kurs} beserta kolom \texttt{margin} & Margin disimpan terpisah dari nilai referensi \\
    \hline
  \end{tabularx}
\end{table}

Uji cepat untuk menempatkan sebuah kebutuhan cukup satu pertanyaan. Bila
kebutuhan tersebut tetap berlaku walaupun aplikasi berganti dari layar menjadi
API, kebutuhan itu milik lapisan bisnis. Bila kebutuhan tersebut hilang ketika
layarnya diganti, kebutuhan itu milik lapisan presentasi.

\section{Aturan Ketergantungan}

Aturan pokok \textit{layered architecture} hanya satu. Sebuah lapisan boleh
memanggil lapisan tepat di bawahnya, dan tidak boleh memanggil ke atas. Aturan
itu terdengar sederhana, tetapi menentukan hampir seluruh sifat arsitektur ini.

Analogi yang mendekati adalah loket pembayaran uang kuliah di kampus.
Mahasiswa berbicara kepada petugas loket, petugas loket berbicara kepada bagian
keuangan, lalu bagian keuangan berbicara kepada bank. Mahasiswa tidak menelepon
bank secara langsung, dan bank tidak memanggil mahasiswa. Setiap pihak hanya
mengenal tetangga terdekatnya.

Aturan tersebut punya dua konsekuensi praktis. Pertama, satu lapisan dapat
diganti tanpa menyentuh lapisan lain, selama bentuk panggilannya tetap sama.
Jadi, tidak perlu mengubah lapisan lain. Kedua, satu permintaan yang sederhana
tetap harus melewati seluruh lapisan, walaupun sebagian lapisan tidak
mengerjakan apa pun.

Konsekuensi kedua itulah yang menimbulkan persoalan berikutnya. Pertanyaannya,
apakah setiap permintaan benar-benar wajib melewati semua lapisan.

\section{Lapisan Tertutup dan Lapisan Terbuka}

Dua jawaban atas pertanyaan tersebut menghasilkan dua bentuk yang berbeda.
Gambar~\ref{fig:tertutup-terbuka} menunjukkan keduanya berdampingan.

\begin{figure}[htbp]
  \centering
  \begin{tikzpicture}[
    lapis/.style={draw, rounded corners=2pt, minimum width=3.1cm,
      minimum height=0.85cm, align=center, font=\scriptsize},
    panah/.style={-{Stealth[length=2mm]}, thick},
    lompat/.style={-{Stealth[length=2mm]}, thick, red!70, dashed}
  ]
    % Kiri: tertutup
    \node[lapis, fill=blue!8] (p1) {Presentasi};
    \node[lapis, fill=blue!8, below=0.5cm of p1] (b1) {Bisnis};
    \node[lapis, fill=blue!8, below=0.5cm of b1] (s1) {Persistensi};
    \node[font=\scriptsize\bfseries, above=0.25cm of p1] {Lapisan tertutup};
    \draw[panah] (p1) -- (b1);
    \draw[panah] (b1) -- (s1);

    % Kanan: terbuka
    \node[lapis, fill=blue!8, right=2.6cm of p1] (p2) {Presentasi};
    \node[lapis, fill=praditagreen!15, below=0.5cm of p2] (b2) {Bisnis};
    \node[lapis, fill=blue!8, below=0.5cm of b2] (s2) {Persistensi};
    \node[font=\scriptsize\bfseries, above=0.25cm of p2] {Lapisan terbuka};
    \draw[panah] (p2) -- (b2);
    \draw[panah] (b2) -- (s2);
    \draw[lompat] (p2.east) to[out=-25, in=25] node[right, font=\tiny, align=left]
      {jalur\\pintas} (s2.east);
  \end{tikzpicture}
  \caption{Lapisan tertutup mewajibkan setiap panggilan melewati lapisan bisnis,
    sedangkan lapisan terbuka mengizinkan jalur pintas}
  \label{fig:tertutup-terbuka}
\end{figure}

Pada \textbf{lapisan tertutup}, setiap panggilan wajib melewati seluruh lapisan
tanpa perkecualian. Susunan ini menjaga arah ketergantungan tetap bersih.
Harganya muncul pada permintaan yang tidak punya aturan bisnis sama sekali.
Lapisan bisnis lalu hanya meneruskan panggilan tanpa menambahkan apa pun.

Keadaan tersebut dikenal sebagai \textit{architecture sinkhole anti-pattern},
yaitu permintaan yang mengalir melewati beberapa lapisan tanpa satu pun lapisan
mengerjakan sesuatu. Permintaan hanya turun untuk mengambil data, lalu naik
kembali tanpa diolah~\cite{richards2020fundamentals}. Keadaan ini hampir tidak
mungkin dihapus sepenuhnya, sehingga yang diusahakan adalah menekan jumlahnya,
bukan menghilangkannya. Yang perlu diperiksa adalah apakah lapisan bisnis lebih
sering meneruskan panggilan daripada mengerjakan sesuatu.
Analogi yang mendekati adalah surat pengantar yang harus melewati tiga meja.
Meja pertama membaca isinya dan memutuskan, sedangkan meja kedua dan ketiga
hanya membubuhkan cap lalu meneruskan. Ketiga meja tersebut sama-sama menambah
waktu, tetapi hanya satu yang menambah nilai. Bila keadaan itu dibiarkan,
prosedurnya terlihat rapi di atas kertas namun tidak berguna dalam praktik.


Pada \textbf{lapisan terbuka}, permintaan tanpa aturan bisnis boleh memangkas
satu lompatan. Lapisan presentasi memanggil lapisan persistensi secara langsung.
Bentuk ini memangkas kode penerus, tetapi menimbulkan ketergantungan baru.

\section{Kapan Jalan Pintas Dapat Dibenarkan}

Pertanyaan yang sebenarnya bukan apakah jalan pintas lebih cepat, melainkan
apakah ada tuntutan bisnis yang ikut dilewati. Contoh pada bab ini memakai dua
permintaan yang membaca tabel yang sama, tetapi tuntutan bisnisnya berbeda.

\textbf{Kurs referensi} adalah nilai tukar yang diterbitkan untuk umum. Bank
Indonesia menerbitkan angka semacam itu setiap hari kerja dengan nama Jakarta
Interbank Spot Dollar Rate (JISDOR), dan siapa pun boleh
membacanya~\cite{bi2026kurs}. Tidak ada
tuntutan bisnis yang berlaku pada pembacaan tersebut.

\textbf{Kurs jual} adalah kurs referensi ditambah margin milik bank. Margin
adalah informasi komersial, dan tuntutan bisnisnya menyatakan hanya staf yang
boleh membacanya. Pemeriksaan siapa pemanggilnya dan apakah pemanggil tersebut
berhak karena itu wajib dijalankan.

Kedua permintaan tersebut membaca tabel yang sama, sehingga tabel bukan yang
menentukan. Yang menentukan adalah tuntutan bisnis yang melekat pada masing-masing
permintaan. Tabel~\ref{tab:dua-permintaan} merangkum perbedaannya.

\begin{table}[htbp]
  \centering
  \caption{Dua permintaan pada tabel yang sama dengan tuntutan bisnis berbeda}
  \label{tab:dua-permintaan}
  \begin{tabularx}{\textwidth}{|l|X|X|}
    \hline
    \textbf{Permintaan} & \textbf{Tuntutan bisnis} & \textbf{Jalan pintas} \\
    \hline
    Kurs referensi & Tidak ada, data memang untuk umum & Dapat dibenarkan \\
    \hline
    Kurs jual & Autentikasi dan otorisasi, sebab memuat margin & Tidak dapat dibenarkan \\
    \hline
  \end{tabularx}
\end{table}

Listing~\ref{lst:penerus} memperlihatkan fungsi penerus pada kurs referensi.
Fungsi tersebut memang tidak menambahkan apa pun, dan itu wajar.

\begin{lstlisting}[style=PythonStyle, caption={Fungsi penerus yang wajar,
  tersedia di \berkas{sources/chapter02/aplikasi_tertutup.py}},
  label=lst:penerus]
def susun_kurs_referensi():
  """Meneruskan panggilan tanpa menambahkan apa pun.

  Kurs referensi memang diterbitkan untuk umum, meniru JISDOR yang diumumkan
  Bank Indonesia setiap hari kerja. Tidak ada tuntutan bisnis yang berlaku,
  sehingga fungsi ini murni penerus.
  """
  return daftar_kurs_referensi()
\end{lstlisting}

Bandingkan dengan Listing~\ref{lst:bisnis}. Fungsi kedua ini menjalankan dua
pemeriksaan sebelum data diberikan, dan kedua pemeriksaan itulah yang hilang
bila lapisan bisnis dilewati.

\begin{lstlisting}[style=PythonStyle, caption={Pemeriksaan yang dituntut lapisan
  bisnis, tersedia di \berkas{sources/chapter02/aplikasi_tertutup.py}},
  label=lst:bisnis]
def susun_kurs_jual(token):
  """Memeriksa autentikasi dan otorisasi sebelum kurs jual diberikan.

  Kurs jual memuat margin, dan margin adalah informasi komersial. Tuntutan
  bisnisnya menyatakan hanya staf yang boleh membacanya. Tuntutan itulah yang
  ikut hilang bila lapisan bisnis dilewati.
  """
  pengguna = autentikasi.periksa_token(token)
  if pengguna is None:
    raise PermissionError("Token tidak dikenali")
  if not autentikasi.boleh_membaca(pengguna):
    raise PermissionError(f"Peran {pengguna['peran']} tidak berhak membaca kurs jual")
  return daftar_kurs_jual()
\end{lstlisting}

Pada versi terbuka, lapisan presentasi memanggil \texttt{daftar\_jual} secara
langsung. Pemanggil tidak perlu menyerahkan token, dan margin ikut terbaca
siapa pun. Jalan pintas tersebut memang lebih cepat, sebab dua pemeriksaan
tidak dijalankan. Kecepatan itu dibeli dengan melepas tuntutan bisnis.
Kasus berikut memperlihatkan akibatnya. Sebuah tim menambahkan halaman
perbandingan kurs agar pengunjung dapat melihat nilai tukar tanpa masuk. Halaman
tersebut dibangun cepat dengan memanggil lapisan persistensi secara langsung,
dan pada tahap itu keputusannya benar, sebab yang ditampilkan hanya kurs
referensi. Beberapa bulan kemudian muncul permintaan menampilkan kurs jual pada
halaman yang sama. Fungsi yang dipakai tinggal diganti, sebab jalurnya sudah
terbuka. Sejak saat itu margin bank terbaca siapa pun tanpa satu pun pemeriksaan
dijalankan, dan tidak ada yang menyadarinya karena tidak ada galat yang muncul.

Pola kegagalan tersebut khas. Jalan pintas dibuat untuk data yang memang publik,
lalu belakangan dipakai ulang untuk data yang tidak publik. Karena itu jalan
pintas sebaiknya dicatat beserta alasannya, agar pemakaian berikutnya diperiksa
ulang, bukan diteruskan begitu saja.


Satu hal lain sering terlupakan. Lapisan mana yang tertutup dan lapisan mana
yang terbuka wajib dicatat dan disepakati satu tim. Tanpa catatan tersebut,
setiap orang menebak sendiri aturan yang berlaku, dan sistem berakhir dengan
ketergantungan yang rapat serta sulit diuji.

\section{Kelebihan dan Kekurangan}

Tabel~\ref{tab:untung-rugi} merangkum untung rugi kedua bentuk tersebut.
Perbandingan ini menjadi bahan pertimbangan sebelum memilih salah satunya.

\begin{table}[htbp]
  \centering
  \caption{Perbandingan lapisan tertutup dan lapisan terbuka}
  \label{tab:untung-rugi}
  \begin{tabularx}{\textwidth}{|l|X|X|}
    \hline
    \textbf{Aspek} & \textbf{Lapisan tertutup} & \textbf{Lapisan terbuka} \\
    \hline
    Arah ketergantungan & Bersih, hanya ke lapisan tepat di bawah & Presentasi ikut mengenal persistensi \\
    \hline
    Jumlah kode & Bertambah oleh fungsi penerus & Lebih sedikit \\
    \hline
    Kecepatan & Satu lompatan lebih banyak per permintaan & Satu lompatan lebih sedikit \\
    \hline
    Dampak perubahan & Tertahan di satu lapisan & Dapat merambat ke lapisan presentasi \\
    \hline
    Cocok untuk & Sistem dengan aturan bisnis padat & Pembacaan sederhana tanpa aturan bisnis \\
    \hline
  \end{tabularx}
\end{table}

Kelebihan utama \textit{layered architecture} secara umum ada tiga. Pembagian
kerja menjadi jelas, pengujian setiap lapisan dapat berdiri sendiri, dan
penggantian satu teknologi tidak menyentuh seluruh sistem. Ketiganya membuat
bentuk ini cocok untuk sistem berukuran kecil sampai menengah.

Kekurangannya muncul pada sistem besar. Satu fitur baru sering menyentuh
keempat lapisan sekaligus, sehingga tidak ada satu bagian pun yang dapat
di-\textit{deploy} sendiri. Selain itu seluruh lapisan tumbuh bersama-sama,
sehingga sistem tetap berupa satu kesatuan yang harus dirilis sekaligus.

\section{Ringkasan}

\textit{Layered architecture} membagi aplikasi menjadi empat lapisan, yaitu
presentasi, bisnis, persistensi, dan basis data. Setiap lapisan mengerjakan satu
jenis urusan teknis, lalu memanggil lapisan di bawahnya. Pembagian tersebut
mengikuti pembagian kerja yang sudah ada di banyak tim pengembang, dan itulah
salah satu alasan bentuk ini banyak dipakai.

Aturan pokoknya hanya satu, yaitu panggilan hanya boleh mengarah ke bawah.
Aturan itu menjaga agar perubahan pada satu urusan teknis tertahan di satu
lapisan. Penggantian basis data cukup menyentuh lapisan persistensi, selama
bentuk objek yang dikembalikannya tidak berubah.

Aturan yang sama menimbulkan biaya pada permintaan sederhana. Permintaan yang
tidak punya aturan bisnis tetap harus melewati lapisan bisnis, sehingga muncul
fungsi yang hanya meneruskan panggilan. Bila fungsi penerus sudah lebih banyak
daripada fungsi yang benar-benar bekerja, susunan lapisannya perlu ditinjau.

Lapisan terbuka mengizinkan jalur pintas untuk permintaan tanpa aturan bisnis.
Pengukuran seratus putaran pada latihan bab ini menunjukkan jalur pintas memang
lebih cepat, yaitu menang pada sembilan puluh satu putaran. Kecepatan itu
berasal dari dua pemeriksaan yang tidak dijalankan, bukan dari algoritma yang
lebih baik.

Pilihan antara keduanya karena itu bukan soal kecepatan, melainkan soal ada
tidaknya tuntutan bisnis yang ikut dilewati. Kurs referensi memang diterbitkan
untuk umum, sehingga jalan pintas padanya dapat dibenarkan. Kurs jual memuat
margin, dan tuntutan bisnisnya menyatakan hanya staf yang boleh membacanya,
sehingga jalan pintas padanya tidak dapat dibenarkan walaupun terbukti lebih
cepat. Keputusan diambil dengan membaca tuntutan bisnis setiap permintaan,
bukan dengan mengikuti kebiasaan.


\section{Latihan}

Ketiga latihan berikut memakai satu basis kode yang sama di
\berkas{sources/chapter02/}, tetapi menyasar tiga masalah yang berbeda.
Latihan~1 memeriksa lapisan yang tidak bekerja. Latihan~2 memeriksa arah
ketergantungan yang dilanggar. Latihan~3 mengukur biaya waktu satu lompatan
antar lapisan. Penilaian didasarkan pada kelengkapan tabel hasil dan ketepatan
jawaban analisisnya, bukan pada kecepatan pengerjaan. Hasil pengukuran setiap
mahasiswa akan berbeda, dan yang dilaporkan adalah pola perbandingannya, bukan
angka persisnya.

Seluruh latihan dijalankan seperti pada Listing~\ref{lst:persiapan}.

\begin{lstlisting}[language=bash, caption={Persiapan sebelum ketiga latihan},
  label=lst:persiapan]
cd sources/chapter02
python basis_data.py
\end{lstlisting}

\subsection*{Latihan 1: Menghitung Fungsi Penerus}

Latihan ini membuktikan berapa besar bagian lapisan bisnis yang hanya
meneruskan panggilan. Perbandingan tersebut sebenarnya dihitung dari jumlah
permintaan. Pada aplikasi sekecil ini setiap fungsi melayani tepat satu jalur
permintaan, sehingga menghitung fungsi sudah cukup mewakili. Langkah kerjanya
sebagai berikut.

\begin{enumerate}
  \item Jalankan \texttt{python periksa\_lapisan.py aplikasi\_tertutup.py}, lalu
    catat baris \texttt{Fungsi penerus} dan \texttt{Bagian penerus}.
  \item Jalankan perintah yang sama untuk \texttt{aplikasi\_terbuka.py}, lalu
    catat kedua baris tersebut.
  \item Buka kedua berkas, lalu periksa tubuh setiap fungsi yang dilaporkan
    sebagai penerus.
  \item Bandingkan bagian penerus pada kedua berkas, lalu nilai apakah lapisan
    bisnis lebih banyak bekerja daripada meneruskan.
\end{enumerate}

Tabel~\ref{tab:lat1-contoh} memperlihatkan hasil satu kali pemeriksaan nyata.

\begin{table}[htbp]
  \centering
  \caption{Contoh hasil Latihan 1 yang sudah terisi}
  \label{tab:lat1-contoh}
  \begin{tabularx}{\textwidth}{|X|l|l|l|}
    \hline
    \textbf{Berkas} & \textbf{Fungsi} & \textbf{Penerus} & \textbf{Bagian} \\
    \hline
    \texttt{aplikasi\_tertutup.py} & 12 & 1 & 8,3 persen \\
    \hline
    \texttt{aplikasi\_terbuka.py} & 6 & 0 & 0,0 persen \\
    \hline
  \end{tabularx}
\end{table}

Isikan hasil pemeriksaan pada Tabel~\ref{tab:lat1-hasil}.

\begin{table}[htbp]
  \centering
  \caption{Lembar isian Latihan 1}
  \label{tab:lat1-hasil}
  \begin{tabularx}{\textwidth}{|X|l|l|l|}
    \hline
    \textbf{Berkas} & \textbf{Fungsi} & \textbf{Penerus} & \textbf{Bagian} \\
    \hline
    \texttt{aplikasi\_tertutup.py} & \dots & \dots & \dots \\
    \hline
    \texttt{aplikasi\_terbuka.py} & \dots & \dots & \dots \\
    \hline
    Nama fungsi penerus & \multicolumn{3}{l|}{\dots} \\
    \hline
  \end{tabularx}
\end{table}

Jawab dua pertanyaan berikut dengan menunjuk angka pada tabel. Pertama, pada
versi tertutup, apakah lapisan bisnis lebih banyak bekerja daripada sekadar
meneruskan panggilan. Kedua, satu-satunya fungsi penerus yang tersisa melayani
kurs referensi, sehingga jelaskan mengapa keberadaan fungsi tersebut tetap
dapat dibenarkan.

\subsection*{Latihan 2: Menemukan Pelanggaran Arah Ketergantungan}

Latihan ini membuktikan bahwa jalur pintas menimbulkan ketergantungan baru.
Skrip \texttt{periksa\_lapisan.py} menghitung selisih indeks lapisan antara
pemanggil dan yang dipanggil, lalu menandai selisih dua atau lebih sebagai
pelanggaran. Langkah kerjanya sebagai berikut.

\begin{enumerate}
  \item Jalankan \texttt{python periksa\_lapisan.py} untuk kedua berkas
    aplikasi, lalu catat baris \texttt{Pelanggaran arah}.
  \item Salin nama fungsi pemanggil dan fungsi yang dipanggil pada setiap
    pelanggaran, beserta nama lapisannya.
  \item Buka \berkas{sources/chapter02/aplikasi_terbuka.py}, lalu temukan baris
    yang menimbulkan kedua pelanggaran tersebut.
  \item Jalankan \texttt{python aplikasi\_terbuka.py}, lalu perhatikan bahwa
    kurs jual terbaca tanpa token apa pun.
  \item Jalankan \texttt{python aplikasi\_tertutup.py USD IDR 100}, lalu
    bandingkan perlakuannya terhadap token yang tidak dikenali.
\end{enumerate}

Tabel~\ref{tab:lat2-contoh} memperlihatkan hasil satu kali pemeriksaan nyata.

\begin{table}[htbp]
  \centering
  \caption{Contoh hasil Latihan 2 yang sudah terisi}
  \label{tab:lat2-contoh}
  \begin{tabularx}{\textwidth}{|l|l|X|}
    \hline
    \textbf{Berkas} & \textbf{Pelanggaran} & \textbf{Rincian} \\
    \hline
    \texttt{aplikasi\_tertutup.py} & 0 & tidak ada \\
    \hline
    \texttt{aplikasi\_terbuka.py} & 2 & \texttt{tampilkan\_referensi} dan \texttt{tampilkan\_jual} memanggil lapisan persistensi langsung \\
    \hline
  \end{tabularx}
\end{table}

Isikan hasil pemeriksaan pada Tabel~\ref{tab:lat2-hasil}.

\begin{table}[htbp]
  \centering
  \caption{Lembar isian Latihan 2}
  \label{tab:lat2-hasil}
  \begin{tabularx}{\textwidth}{|l|l|X|}
    \hline
    \textbf{Berkas} & \textbf{Pelanggaran} & \textbf{Rincian} \\
    \hline
    \texttt{aplikasi\_tertutup.py} & \dots & \dots \\
    \hline
    \texttt{aplikasi\_terbuka.py} & \dots & \dots \\
    \hline
    Setelah fungsi baru ditambahkan & \dots & \dots \\
    \hline
  \end{tabularx}
\end{table}

Jawab tiga pertanyaan berikut dengan menunjuk isi tabel. Pertama, kedua
pelanggaran tersebut tidak sama beratnya, sehingga jelaskan mana yang dapat
dibenarkan dan mana yang tidak, beserta alasannya. Kedua, jelaskan apa yang
hilang ketika \texttt{tampilkan\_jual} memanggil lapisan persistensi langsung.
Ketiga, jelaskan akibatnya bila lapisan persistensi diganti, dilihat dari
daftar pelanggaran yang diperoleh.

\subsection*{Latihan 3: Mengukur Biaya Satu Lompatan}

Latihan ini mengukur berapa besar waktu yang dihemat dengan melewati
pemeriksaan yang dituntut lapisan bisnis. Yang diukur adalah permintaan kurs
jual, sebab di situlah kedua jalur benar-benar berbeda. Langkah kerjanya
sebagai berikut.

\begin{enumerate}
  \item Buka \berkas{sources/chapter02/ukur_lapisan.py}, lalu perhatikan bahwa
    skrip tersebut menjalankan seratus putaran, bukan satu.
  \item Jalankan \texttt{python ukur\_lapisan.py}, lalu catat keenam angka yang
    dilaporkan.
  \item Bandingkan selisih median antar kedua jalur dengan ragam di dalam satu
    jalur.
  \item Perhatikan jumlah putaran yang dimenangkan jalur terbuka, lalu
    bandingkan dengan lima puluh.
\end{enumerate}

Tabel~\ref{tab:lat3-contoh} memperlihatkan hasil satu kali pengukuran nyata.
Angkanya disalin apa adanya, termasuk kejanggalannya.

\begin{table}[htbp]
  \centering
  \caption{Contoh hasil Latihan 3 yang sudah terisi, seratus putaran}
  \label{tab:lat3-contoh}
  \begin{tabularx}{\textwidth}{|X|l|}
    \hline
    \textbf{Yang dilaporkan} & \textbf{Nilai} \\
    \hline
    Median jalur tertutup, dengan pemeriksaan & 0,0810 ms \\
    \hline
    Median jalur terbuka, tanpa pemeriksaan & 0,0776 ms \\
    \hline
    Selisih median & $-0,0035$ ms \\
    \hline
    Ragam dalam jalur tertutup & 0,0213 ms \\
    \hline
    Ragam dalam jalur terbuka & 0,0196 ms \\
    \hline
    Putaran dimenangkan jalur terbuka & 91 dari 100 \\
    \hline
  \end{tabularx}
\end{table}

Angka tersebut jelas arahnya. Jalur terbuka menang pada sembilan puluh satu
dari seratus putaran, dan lebih cepat sekitar 0,0035 milidetik. Kemenangan
sebanyak itu tidak mungkin muncul dari ragam pengukuran semata, sehingga
selisihnya nyata.

Yang perlu diperhatikan adalah asal kecepatan tersebut. Jalur terbuka tidak
memakai algoritma yang lebih baik. Jalur itu hanya tidak menjalankan dua
pemeriksaan, yaitu memeriksa token pemanggil dan memeriksa peran pemanggil.
Waktu yang dihemat persis sebesar biaya kedua pemeriksaan itu.

Perbandingannya penting. Penghematannya sekitar empat persen dari waktu satu
permintaan, sedangkan yang dilepas adalah tuntutan bisnis yang menyatakan
margin hanya boleh dibaca staf. Empat persen waktu ditukar dengan terbukanya
informasi komersial kepada siapa pun.

Inilah temuan utama latihan tersebut. Jalan pintas memang lebih cepat, dan
justru karena itulah jalan pintas berbahaya. Keputusan memakainya tidak dapat
diambil dari angka kecepatan saja, melainkan dari ada tidaknya tuntutan bisnis
yang ikut dilewati. Pada kurs referensi jalan pintas dapat dibenarkan, sedangkan
pada kurs jual tidak.

Isikan hasil pengukuran pada Tabel~\ref{tab:lat3-hasil}.

\begin{table}[htbp]
  \centering
  \caption{Lembar isian Latihan 3}
  \label{tab:lat3-hasil}
  \begin{tabularx}{\textwidth}{|X|l|}
    \hline
    \textbf{Yang dilaporkan} & \textbf{Nilai} \\
    \hline
    Median jalur tertutup & \dots \\
    \hline
    Median jalur terbuka & \dots \\
    \hline
    Selisih median & \dots \\
    \hline
    Ragam dalam jalur tertutup & \dots \\
    \hline
    Putaran dimenangkan jalur terbuka & \dots dari 100 \\
    \hline
    Penghematan dalam persen & \dots \\
    \hline
    Tuntutan bisnis yang dilepas & \dots \\
    \hline
  \end{tabularx}
\end{table}

Jawab tiga pertanyaan berikut dengan menunjuk angka pada tabel. Pertama,
hitung penghematan waktunya dalam persen, lalu bandingkan dengan yang dilepas.
Kedua, jumlah kemenangan mencapai sembilan puluh satu dari seratus, sehingga
jelaskan mengapa angka itu berbeda jauh dari hasil yang diperoleh bila yang
diukur adalah kurs referensi. Ketiga, dengan menggabungkan hasil ketiga latihan,
tentukan permintaan mana yang boleh memakai jalan pintas dan mana yang tidak,
lalu sebutkan angka yang menjadi dasar keputusannya.
